\chapter{Python/Dask calculations, diagnostics and independent checks}
\label{sec:computational-verification}

\section{Maintained implementation and mathematical coverage}
The active numerical implementation uses Python, Dask, NumPy and SciPy;
the model and mathematical proofs are independent of the chosen
programming language. The common Dask graph generates Gaussian
whole-population and conditional score vectors, computes the sample
moments from seeded, chunked arrays, and writes reproducible tables.
Independent Python reference calculations use the standard library's
rational arithmetic and established SciPy/NumPy numerical methods
rather than a second unreviewed probability implementation.

The nine IQ-150 Monte Carlo examples, finite covariance endpoints and
analytically derived limits are compared against explicit formulas.
The main mathematical theorem additionally covers every integer
\(n\geq1\) of identical starting observations \(a\). An exact rational
Python calculator for arbitrary \(n,k\geq1\) is maintained in
\texttt{library/python/sequential.py}; its independent regression tests
check conditional means and variances using matrix Schur complements,
including \(n=5\) and singular correlation endpoints. These are exact
algebraic and finite-dimensional checks, not new Dask replications.
A pass on the recorded \(n\in\{1,2,3\}\) grid only checks those numerical
instances and cannot establish a uniform sampling-error bound for all \(n\).
Floating-point rounding, finite-state pseudorandom generators and
error bounds for numerical integration require separate evaluation.

\section{What the Python/Dask simulations sample}
For each modeled finite case the population actual-mean law has mean
zero and variance \(d_N/N\), whereas the conditional actual-mean law
has mean \(B_{r,k}\) and variance \(k(1-c)d_N/(N^2d_r)\).
Their IQ-score centers are \(100\) and \(100+15B_{r,k}\);
limits and singular cases are treated according to their stated
mathematical domains. The active experiment constructs complete
finite score vectors for the nonnegative-correlation configurations
reported here; it does not claim a second native scalar experiment.
A zero-variance case is deterministic and must not divide by zero.

For a prescribed sample size \(n\), an unbiased sample variance uses
denominator \(n-1\). The numerical checks compare standardized sample
means and positive-variance chi-square pivots with established SciPy
references. These diagnostics are *conditional on independently sampled
Gaussian model variables*; a finite pseudorandom realization is not a
formal proof of that premise. The propositions below describe the
ideal mathematical sampling laws, not a Rocq certification of Python.


\begin{proposition}[Reference laws for the simulation diagnostics]
\label{prop:simulation-reference-laws}
For an integer \(n\geq2\), if \(Y_1,\ldots,Y_n\) are independent \(N(m,v)\) variables, \(v>0\),
then \(\sqrt{n}(\overline Y-m)/\sqrt v\) is standard normal, and
\((n-1)S^2/v\) has the chi-square distribution with \(n-1\) degrees
of freedom, where \(S^2=\sum_i(Y_i-\overline Y)^2/(n-1)\).
\end{proposition}
\begin{proof}
Let \(Z_i=(Y_i-m)/\sqrt v\). Extend the unit vector
\(n^{-1/2}\one_n\) to an orthonormal basis, and let \(Q\) have those
basis vectors as rows, with that vector first. Such an extension follows
by finite-dimensional orthogonal decomposition (or Gram--Schmidt), as in
\citet[Theorem 6.36, p.~203]{axler2024}. By Theorem~\ref{thm:gaussian}, \(T=QZ\)
is Gaussian with covariance \(QQ^{\T}=I\); its components are therefore
independent standard normals. Its first component is
\(T_1=\sqrt n(\overline Y-m)/\sqrt v\). Preservation of squared norm gives
\[
 \frac{1}{v}\sum_i(Y_i-\overline Y)^2
 =\sum_i Z_i^2-T_1^2=\sum_{j=2}^{n}T_j^2.
\]
The chi-square distribution with \(n-1\) degrees of freedom is, by
definition, the law of the sum of squares of \(n-1\) independent standard
normals. This proves both claims from Gaussian closure and orthogonality;
SciPy evaluates the corresponding normal and chi-square quantiles, rather
than supplying a new mathematical assumption.
\end{proof}

\begin{corollary}[Inversion and finite-grid diagnostic coverage]
\label{cor:simulation-diagnostic-coverage}
Let \(0<\alpha<1\), and let \(q_u\) denote a \(\chi^2_{n-1}\) quantile.
Under the hypotheses
of Proposition~\ref{prop:simulation-reference-laws}, the reference band
\([vq_{\alpha/2}/(n-1),vq_{1-\alpha/2}/(n-1)]\) contains \(S^2\)
with probability \(1-\alpha\). An exact confidence interval for the
unknown SD is instead
\[
 \left[\sqrt{\frac{(n-1)S^2}{q_{1-\alpha/2}}},
       \sqrt{\frac{(n-1)S^2}{q_{\alpha/2}}}\right].
\]
For a fixed finite collection of diagnostics having individual false-alarm
probabilities at most \(\alpha_j\), the probability of any false alarm
is at most \(\sum_j\alpha_j\), even when diagnostics share random draws.
\end{corollary}
\begin{proof}
By the continuous pivot law,
\(P(q_{\alpha/2}\leq(n-1)S^2/v\leq q_{1-\alpha/2})=1-\alpha\).
Solving for \(S^2\) gives the reference band; solving for positive \(v\)
and taking the increasing square root gives the SD interval. This is the
standard pivotal inversion argument, derived from the proved sampling law.
For false-alarm events \(F_j\), the pointwise inequality
\(1_{\cup_jF_j}\leq\sum_j1_{F_j}\), integrated by linearity, gives
\(P(\cup_jF_j)\leq\sum_jP(F_j)\leq\sum_j\alpha_j\). This is
subadditivity of probability \citep[Section 1.1]{durrett2019}. Counting
two tails separately is equivalent to adding their two probabilities.
The bound justifies conservative choices for a prespecified finite grid;
it neither proves PRNG independence nor certifies arbitrary future grids.
Zero-variance cases use constant output and do not divide by \(v\).
\end{proof}

\section{Deal contracts, reproducibility and verification boundaries}
The maintained Python implementation declares value-domain,
shape, exception, frame, and side-effect properties through Deal.
Its contracts, static lint, property tests, source-order report
checks and bounded symbolic checks are different kinds of
evidence; none alone certifies SciPy or NumPy numerical internals.
The separately maintained Rocq library checks only its formal
propositions, under their explicit real and finite-list hypotheses.

The Dask configuration, seeded replicates and generated tables
remain the numerical provenance for the six-volume project.
The discontinued alternative implementation is not part of the
current build or verification gate. Its older claims are preserved
only by Git history, not carried forward as current-source proofs.
No fresh experimental run is asserted by the present editorial
migration. Real-world IQ-test validity, participant independence,
measurement reliability and the Gaussian population model remain
empirical assumptions rather than program-verification conclusions.

